% GreenAI technical documentation v2. Compile from the repository root:
% tectonic -b https://data1.fullyjustified.net/tlextras-2022.0r0.tar \
%   -o output/pdf docs/greenai_v2.tex
\documentclass[11pt,a4paper]{article}
\usepackage{fontspec}
\usepackage{polyglossia}
\setmainlanguage{russian}
\setotherlanguage{english}
\setmainfont{Arial}
\setsansfont{Arial}
\setmonofont{Consolas}
\usepackage[margin=19mm,headheight=16pt]{geometry}
\usepackage{amsmath,amssymb}
\usepackage{graphicx}
\usepackage{booktabs,longtable,array}
\usepackage{enumitem}
\usepackage{xcolor}
\usepackage{fancyhdr}
\usepackage{hyperref}
\usepackage{microtype}
\usepackage{float}
\definecolor{ink}{HTML}{183049}
\definecolor{blue}{HTML}{166F9D}
\definecolor{pale}{HTML}{EAF3F8}
\definecolor{green}{HTML}{187357}
\definecolor{muted}{HTML}{566879}
\hypersetup{colorlinks=true,linkcolor=blue,urlcolor=blue,pdftitle={GreenAI: архитектура и алгоритмы, версия 2},pdfauthor={Команда GreenAI}}
\setlength{\parindent}{0pt}
\setlength{\parskip}{5pt}
\setlist[itemize]{leftmargin=5mm,itemsep=2pt,topsep=3pt}
\setlist[enumerate]{leftmargin=6mm,itemsep=2pt,topsep=3pt}
\pagestyle{fancy}
\fancyhf{}
\fancyhead[L]{\small\color{muted}GreenAI / техническая документация}
\fancyhead[R]{\small\color{muted}v2 · 27.09.2026}
\fancyfoot[C]{\small\color{muted}\thepage}
\renewcommand{\headrulewidth}{0.4pt}
\newcommand{\code}[1]{\texttt{#1}}
\newcommand{\src}[1]{\par{\small\color{muted}Реализация: \texttt{#1}.}\par}
\newcommand{\alg}[1]{\medskip\noindent\colorbox{pale}{\parbox{\dimexpr\linewidth-2\fboxsep}{\textbf{#1}}}\smallskip}
\newcommand{\fig}[3]{\begin{figure}[H]\centering\includegraphics[width=#2\linewidth]{#1}\caption{#3}\end{figure}}
\begin{document}

\begin{titlepage}
\vspace*{28mm}
{\color{blue}\rule{35mm}{2.5pt}}\par\vspace{13mm}
{\Huge\bfseries\color{ink}GreenAI}\par\vspace{5mm}
{\LARGE Архитектура и алгоритмы}\par\vspace{3mm}
{\Large Техническая документация · версия 2}\par
\vspace{17mm}
\begin{minipage}{0.84\linewidth}
Автоматическое формирование плана озеленения по топографической подоснове DXF: распознавание объектов, восстановление геометрии, расчёт ограничений, размещение растений и объяснение решений.
\end{minipage}
\vfill
\noindent\colorbox{pale}{\parbox{0.82\linewidth}{\textbf{Статус документа.} Формулы ниже описывают действующую реализацию в репозитории на 27 сентября 2026 года. Параметры правил и моделей читаются из конфигурации и базы данных; приведённые численные примеры служат для пояснения, а не заменяют запись конкретного прогона.}}\par
\vspace{12mm}
\noindent Москва · 2026
\end{titlepage}

\tableofcontents
\clearpage

\section{Назначение и границы решения}
GreenAI получает топографический DXF, выделяет рабочую территорию, пригодные поверхности, дороги, здания и инженерные сети, затем строит план деревьев, кустарниковых массивов и травянистого покрытия. Результат включает новый DXF с посадками, геометрические слои, отчёты о проверках и объяснения для каждой принятой посадки. Исходный чертёж сохраняется отдельно.

\textbf{В MVP входят:} извлечение DXF и блоков, проверка единиц, нормализация геометрии, четыре ONNX-классификатора сетей (вода, газ, теплосеть, силовой кабель), геометрическое восстановление допустимых разрывов, построение базовой и видовых зон, подбор растений, планирование посадок, DXF/PDF-выгрузка, диагностика и независимая проверка результата. Основной интерфейс --- CLI и скрипты; развёртывание --- Docker Compose.

\textbf{Следующие релизы:} расширение проверок на все 20 улиц, увеличение размеченной выборки по сетям, интерактивная правка плана с мгновенной перепроверкой, полные правила для канализации, водостока и связи, устойчивый импорт нестандартных DWG/XREF и геопривязка по подтверждённым опорным точкам. Воздушные ЛЭП со стрелочными обозначениями пока требуют проверки исходных условных знаков и трассы.

\section{Модель данных и поток обработки}
Единица учёта после извлечения --- объект с типом, исходным слоем, геометрией и идентификатором. Промежуточные объекты хранятся в JSONL/GeoJSONL. База PostGIS содержит каталог растений, правила и нормативные ссылки. Для конкретного прогона фиксируются входной DXF, карта слоёв, масштаб единиц, выбранная модель, настройки планировщика и отчёты.

\fig{diagrams/system_design.png}{0.98}{Компоненты и обмен данными. Подписи на схеме описывают роль каждого блока в расчёте.}

\small\begin{longtable}{p{0.22\linewidth}p{0.39\linewidth}p{0.30\linewidth}}
\toprule
Стадия & Вход и действие & Выход / контроль\\\midrule\endhead
Единицы и DXF & Чтение \code{\$INSUNITS}, извлечение modelspace и вложенных блоков, сопоставление слоёв & \code{dxf\_units\_report.json}, \code{extracted\_objects.jsonl}\\
Подготовка & Нормализация поверхностей, дорог и препятствий; классификация и восстановление сетей & \code{normalized\_objects.geojsonl}, \code{cleaned\_utilities.geojsonl}, \code{reconstructed\_\allowbreak utilities.geojsonl}\\
Ограничения & Физические исключения, буферы правил, контроль покрытия исходными данными & \code{constraint\_map.geojsonl}, \code{plant\_allow\_zones.geojsonl}\\
План & Выбор вида, композиции, кандидатов и их проверка & \code{planting\_plan.geojsonl}, \code{planting\_decisions.geojsonl}\\
Выпуск & Добавление слоёв к копии чертежа, независимая верификация & \code{result\_with\_\allowbreak planting\_plan.dxf}, \code{verification\_report.json}, PDF\\
\bottomrule
\end{longtable}
\normalsize

\fig{diagrams/process_flow.png}{0.98}{Путь от исходного DXF к проверенному чертежу и дополнительной визуализации.}

\section{Единицы измерения и геометрическая основа}
Обозначим через $q$ число единиц DXF в одном метре. Для расстояния правила $d_m$ и площади $A_{DXF}$ применяется
\begin{equation}
d_{DXF}=q\,d_m,\qquad A_{m^2}=\frac{A_{DXF}}{q^2}.
\end{equation}
Масштаб определяется из заголовка DXF или явно подтверждённого параметра запуска. При неизвестных единицах расчёт останавливается: отступы в метрах без масштаба нельзя достоверно перенести в чертёж.

Пусть $W$ --- полигон границы работ, $S$ --- объединение подтверждённых пригодных поверхностей. Стартовая поверхность $C$ определяется так:
\begin{equation}
C=\begin{cases}W\cap S,& S\ne\varnothing,\\W,& S=\varnothing\ \text{(резервный режим с предупреждением).}\end{cases}
\end{equation}
Пусть $E$ --- объединение физических препятствий: дорог, тротуаров, твёрдых покрытий, зданий, пятен колодцев и камер теплосети. \textbf{Базовая зона} $B$ --- пространство, в котором после вычитания этих объектов остаётся поверхность для дальнейшего анализа:
\begin{equation}
B=\operatorname{Polygonal}(C\setminus E).
\end{equation}
Это ещё не разрешение на конкретное дерево или кустарник: сетевые и объектные отступы применяются на следующем шаге. В резервном режиме площадь $W$ не считается подтверждённой пригодной поверхностью; статус и диагностические проверки сохраняют эту неопределённость.
\src{src/geometry/constraint\_builder.py; src/geometry/sidewalk\_geometry.py}

\section{Распознавание и восстановление инженерных сетей}
\subsection{Классификация графических примитивов}
В слое сети могут находиться оси труб, стрелки, выноски и подписи. Для каждого линейного примитива строится 19-мерный вектор признаков $x_i$: длина, прямолинейность, габариты, число вершин, степень связности концов, продолжения, соседние штрихи и характеристики связной компоненты. Отдельная модель $M_t$ для типа сети $t$ возвращает $p_i=M_t(x_i)$.

\begin{equation}
\operatorname{decision}(p_i)=
\begin{cases}
\text{accepted}, & p_i\ge \tau_{a,t},\\
\text{manual\_review}, & \tau_{r,t}\le p_i<\tau_{a,t},\\
\text{rejected}, & p_i<\tau_{r,t}.
\end{cases}
\end{equation}
Пороги $\tau_{a,t}$ и $\tau_{r,t}$ хранятся вместе с ONNX-моделью в JSON-метаданных; вектор подаётся в заданном там порядке как \code{float32}. При обучении пространственные группы позволяют отделять участки проверки от обучающих. Для кабеля порог принятия ориентирован на точность, для остальных сетей --- на полноту. Валидационные показатели относятся к размеченным выборкам и не гарантируют такую же точность на новой улице.

\alg{Алгоритм 1. Очистка сети}
\begin{enumerate}
\item Извлечь из целевых слоёв линейные примитивы и построить их признаки с учётом соседей.
\item Выполнить ONNX-инференс отдельно для воды, газа, теплосети и кабеля.
\item Разделить примитивы по двум порогам; сохранить принятые оси и отдельный слой ручной проверки.
\item Передать в восстановление только принятую геометрию и сохранить исходные ID для трассировки.
\end{enumerate}
\src{src/detection/utilities/detector.py;\\ src/detection/utilities/onnx\_model.py; models/utility\_detector/latest/metadata.json}

\subsection{Короткие разрывы и неоднозначность}
Открытые концы принятых линий индексируются пространственно. Для кандидата продолжения длиной $g$ вычисляются два угла $\alpha_1,\alpha_2$ между направлением линии у концов и соединяющим вектором. Кандидат проходит геометрический фильтр $g\le g_{\max}$ и $\alpha_1,\alpha_2\le\alpha_{\max}$. Его базовый скоринг:
\begin{align}
D&=1-\frac{\alpha_1+\alpha_2}{2\alpha_{\max}}, & L&=1-\frac{g}{g_{\max}},\\
s_{\mathrm{cont}}&=0.7D+0.3L.&&
\end{align}
Для присоединения конца к внутренней точке соседней линии используется проекция и весовая схема $s_{\mathrm{junction}}=0.65D+0.35L$; в этом случае $D$ зависит от одного угла. Настройки содержат лимиты разрыва, углов и пороги решений. При параллельных трассах поддержка соседним сходным продолжением повышает оценку. Если несколько кандидатов конкурируют, алгоритм фиксирует неоднозначность и может оставить ветвление по настроенной политике. Принятая соединительная линия записывается отдельно от исходной оси, что позволяет увидеть именно реконструкцию. Для силового кабеля автоматическое соединение разрывов отключено.

\alg{Алгоритм 2. Восстановление фрагмента сети}
\begin{enumerate}
\item Найти свободные концы принятой сети и локальные кандидаты через пространственный индекс.
\item Отсечь дальние и геометрически несовместимые пары; оценить продолжения и присоединения.
\item Учесть параллельность, конкурирующие соединения и заданные пороги.
\item Записать принятые соединители, сомнительные кандидаты и причины решения в отдельные артефакты.
\end{enumerate}
\src{src/detection/network\_reconstructor.py;\\ src/core/network\_reconstructor\_config.yaml}

\section{Правила и допустимые зоны}
Для типа посадки $t\in\{\mathrm{tree},\mathrm{shrub}\}$ активные правила $R_t$ задают целевой объект $G_r$, отступ $d_r$ в метрах и ссылку на НПА. Геометрия запрета для правила и результирующая \textbf{допустимая зона} записываются как
\begin{equation}
X_r=\operatorname{Buffer}(G_r,q d_r),\qquad
A_t=\operatorname{Polygonal}\left(B\setminus\bigcup_{r\in R_t}X_r\right).
\end{equation}
Таким образом, $B$ отвечает на вопрос «где в принципе есть поверхность», а $A_t$ --- «где по доступным данным разрешён центр посадки данного типа». Условия размещения кроны, расстояния между растениями и выбранного вида проверяются затем при генерации плана. Для травянистого покрытия отдельная видовая зона правил не строится: оно занимает оставшуюся часть $B$ с учётом существующих массивов и уже размещённых элементов.

В текущем каталоге заданы следующие минимальные отступы от оси ствола или центра посадки до объекта. Их источник в каталоге --- \href{https://protect.gost.ru/sp/details/f6917ab4-63d8-4ecb-9794-0b4990ba3b99}{СП 42.13330.2026, таблица 6.3}. Для конкретного прогона действует содержимое базы после выполнения \code{seed.py}.
\begin{center}\small
\begin{tabular}{p{0.37\linewidth}ccp{0.30\linewidth}}
\toprule
Объект & Дерево, м & Кустарник, м & Код правила (пример)\\\midrule
Здание & 5 & 1.5 & TREE\_BUILDING\_5\\
Тротуар & 0.7 & 0.5 & TREE\_SIDEWALK\_0\_7\\
Край проезжей части & 2 & 1 & TREE\_ROAD\_EDGE\_2\\
Газопровод & 1.5 & --- & TREE\_GAS\_1\_5\\
Теплосеть & 2 & 1 & TREE\_HEAT\_2\\
Водопровод & 2 & --- & TREE\_WATER\_2\\
Силовой кабель & 2 & 0.75 & TREE\_POWER\_CABLE\_2\\
\bottomrule
\end{tabular}\end{center}
Прочерк означает отсутствие отдельного правила данного типа в текущем каталоге; он не является утверждением о нормативной допустимости посадки.

Если объект правила отсутствует, правило получает статус \code{unavailable}; вычитание по нему не выполняется, а неполнота данных попадает в отчёт. Если сеть имеет сомнительный источник или неполную геометрию, статус зоны указывает на необходимость проверки. Проверка результата отличает вычисленное отсутствие нарушения от отсутствия данных о препятствии.

\fig{diagrams/allow_zone_algorithm.png}{0.96}{Базовая область, буферы правил и зоны для разных типов посадок.}

Пример: при $q=1$ и правиле «кустарник --- теплосеть: 1 м» каждый сегмент подтверждённой теплотрассы расширяется буфером шириной 1 м, после чего буфер вычитается из $B$. Фактическое значение, код правила и ссылка на пункт нормативного документа берутся из записи PostGIS, а не из этой иллюстрации. Для камер теплосети действует самостоятельное физическое исключение контура.
\src{src/rules/plant\_allow\_zone.py; scripts/seed.py}

\section{Планировщик посадок}
Планировщик читает допустимые зоны, каталог морозостойких растений и запрос на композицию. Профиль вида содержит тип, площадь/радиус пятна, шаг между растениями, расстояние до других посадок и причины выбора. Деревья получают точечные центры; кустарник может описываться непрерывным полигоном массива; травянистое покрытие --- полигоном остатка. Выбор вида и композиции фиксируется в объяснении.

Для двух профилей $i,j$ минимальный шаг центров в метрах равен
\begin{equation}
s_{ij}=\begin{cases}
\max(s_i,s_j,r_i+r_j),&t_i=t_j,\\
\max(a_i,a_j,r_i+r_j),&t_i\ne t_j,
\end{cases}
\end{equation}
где $s_i$ --- шаг вида, $a_i$ --- избегаемое расстояние до другого типа, $r_i$ --- радиус пятна. Перед принятием точки проверяется её нахождение в безопасной области с учётом радиуса, существующих деревьев и уже выбранных посадок.

Для площадного заполнения используется сдвинутая по соседним строкам сетка. При шаге $s$ расстояние между строками $h=\sqrt{3}s/2$; каждая вторая строка сдвигается на $s/2$. Для вытянутых фрагментов сначала оценивается продольная ось, затем несколько фаз ряда. В линейном режиме допустимы схемы с числом принятых точек не менее $\lceil0.9N_{\max}\rceil$; из них выбирается композиция с большим числом соседних пар рядов, меньшим числом одиночных рядов и лучшим балансом краёв. Это даёт ровную посадку без существенной потери вместимости.

\alg{Алгоритм 3. Размещение и протокол решений}
\begin{enumerate}
\item Выбрать вид и тип композиции для каждой части зоны; сначала обрабатывать более крупные пятна.
\item Сгенерировать фазы точечной сетки или линейных рядов; проверить каждую точку по зоне, соседям и существующей растительности.
\item Записать принятые точки в план, отклонённые --- в диагностику с фактическим расстоянием и требуемым пределом.
\item Сформировать кустарниковые полигоны и остаточное травянистое покрытие. Для массива оценить количество кустов как $\left\lceil A/(s^2\sqrt{3}/2)\right\rceil$ при положительном $s$.
\item Сохранить выбранную фазу, число точек и список причин принятия или отклонения для последующей проверки.
\end{enumerate}
\src{src/planting/placement\_generator.py; src/planting/service.py}

\fig{diagrams/planting_algorithm.png}{0.97}{Построение композиций и проверка кандидатов.}

\section{Интерпретируемость и независимая проверка}
Каждая принятая посадка получает ID, геометрию, выбранный вид, источник зоны, список пройденных проверок, фактические и требуемые расстояния, код правила и нормативную ссылку с пунктом. Для отклонённого кандидата сохраняется проверка, вызвавшая отказ, и измеренное значение. Подробности распределены по трём файлам:
\begin{itemize}
\item \code{planting\_plan.geojsonl} --- принятый план;
\item \code{planting\_decisions.geojsonl} --- решения по кандидатам;
\item \code{planting\_explanations.md} --- читаемые пояснения.
\end{itemize}
PDF даёт краткое представление для эксперта. Ссылка на НПА берётся из версии правила в базе, поэтому при изменении правил пересчёт создаёт новый набор объяснений.

Интерпретация решения $c$ может быть представлена записью
\begin{equation}
E(c)=\{\operatorname{ID},\operatorname{geometry},\operatorname{species},
\operatorname{checks},\operatorname{measurements},\operatorname{rule\_ids},
\operatorname{references},\operatorname{source\_quality}\}.
\end{equation}
Это структура отчёта, а не дополнительная формула принятия: итоговый статус определяется последовательностью проверок и записывается вместе с ними. Для экспертной оценки важно открывать полный протокол, если краткий PDF скрывает часть кандидатов.

\alg{Алгоритм 4. Контроль готового результата}
\begin{enumerate}
\item Повторно считать геометрию плана и допустимых зон из сохранённых артефактов.
\item Проверить принадлежность центров зонам, непересечение с физическими препятствиями и минимальные расстояния между посадками.
\item Сопоставить исходные объекты с копией DXF и проверить структуру выходного файла.
\item Выдать \code{passed} только для набора выполненных вычислимых проверок; недостаток входных данных оставить явным в отчёте.
\end{enumerate}
\src{scripts/verify\_outputs.py; src/planting/service.py}

\section{Развёртывание и воспроизведение}
\fig{diagrams/deployment.png}{0.91}{Развёртывание: контейнер планировщика, PostGIS и каталог входных/выходных данных.}

На Linux или МосТех.ОС с Docker Engine и Compose из корня репозитория выполнить:
\begin{verbatim}
mkdir -p docker-data
cp /path/to/input.dxf docker-data/input.dxf
docker compose build planner seed
docker compose up -d --wait postgis
docker compose --profile tools run --rm seed
docker compose --profile pipeline run --rm planner \
  /data/input.dxf /data/output
\end{verbatim}
Затем открыть отчёт \code{verification\_report.json} в \code{docker-data/output}. При статусе \code{passed} следует просмотреть \code{result\_with\_planting\_plan.dxf} и PDF-отчёт в CAD и просмотрщике.
Проверка на целевой сборке МосТех.ОС после последнего рефакторинга ещё должна быть зафиксирована отдельным прогоном. На Windows установку выполняет \code{scripts/install.ps1}, а полный запуск --- \code{scripts/run\_pipeline.ps1} с \code{-InputDxf}, \code{-OutputDirectory} и \code{-PipelineMode full}. Повторный запуск в новой выходной папке сохраняет результаты предыдущего расчёта.

HTTP-бэкенда у MVP нет; способы взаимодействия --- CLI, конфигурационные файлы и выходные артефакты. OpenAPI/Swagger потребуется при появлении HTTP API.

\section{Методы, библиотеки и модели}
\begin{longtable}{p{0.26\linewidth}>{\raggedright\arraybackslash}p{0.62\linewidth}}
\toprule
Компонент & Назначение и источник\\\midrule\endhead
Go-парсер DXF & Потоковое извлечение объектов и блоков; \href{https://go.dev/}{Go}, \href{https://pkg.go.dev/gopkg.in/yaml.v3}{yaml.v3}.\\
\href{https://ezdxf.readthedocs.io/}{ezdxf} & Анализ DXF, нормализация и создание выходного чертежа.\\
\href{https://shapely.readthedocs.io/}{Shapely / GEOS} & Буферы, пересечения, разности, пространственные индексы, полигонизация.\\
\href{https://postgis.net/documentation/}{PostGIS} & Хранение правил, видов растений и нормативных ссылок.\\
\href{https://scikit-learn.org/stable/}{scikit-learn} & Обучение четырёх бинарных RandomForest-классификаторов.\\
\href{https://onnx.ai/}{ONNX} / \href{https://onnxruntime.ai/docs/}{ONNX Runtime} & Сериализация моделей и инференс в конвейере. Метаданные с порогами и порядком признаков хранятся рядом с ONNX.\\
\href{https://www.reportlab.com/docs/reportlab-userguide.pdf}{ReportLab} & Формирование PDF-отчёта по посадкам; \href{https://www.blender.org/}{Blender} --- дополнительные изображения до/после.\\
\bottomrule
\end{longtable}

\section{Ограничения и контроль качества входа}
Чертежи улиц пока не имеют единого стандарта слоёв и условных обозначений. Часть объектов находится во вложенных блоках, отдельные DXF возникают из DWG с отсутствующими внешними ссылками или ошибками таблиц CAD. Поэтому карта слоёв и отчёт извлечения должны проверяться на новой улице; ошибочный или пустой источник нельзя считать доказательством свободной территории.

Неполные сведения о подземной сети, воздушной ЛЭП или напряжении требуют отдельной проверки. Стрелочные знаки ЛЭП сами по себе не задают однозначную трассу и охранную зону. Геопривязка локальной системы к WGS84 невозможна без достоверных контрольных точек и определения исходной системы координат. Автоматическая проверка покрывает только описанные и доступные на чертеже ограничения; итоговый проект подлежит инженерной экспертизе.

При переносе на новую улицу минимальная проверка включает: отчёт единиц, число объектов по ключевым типам, положение слоёв относительно границы работ, диагностические слои очищенных сетей и исключений, отчёт зон, протокол отклонённых кандидатов, открытие DXF в целевом CAD и результат \code{verification\_report.json}.

\vfill
{\small\color{muted}Исходный код и настройки являются источником точных параметров конкретного прогона. Документ объясняет метод и точки проверки; численные значения действующих правил определяются сохранённой версией каталога PostGIS и метаданными модели.}
\end{document}
